\documentclass[12pt]{article}
\usepackage{amsmath,amssymb,amsfonts,mathrsfs}
\usepackage[margin=1in]{geometry}

\begin{document}

\begin{center}
{\Large\bfseries Chapter 10, Problem 19}
\end{center}

\noindent\textbf{Problem.} Suppose that we have a linear operator with cubic symmetry, and we apply a small perturbation having the symmetry of a rectangular parallelepiped. Making use of the results of Problem 10.16 and the work of Section 10.8, find how the various possible degenerate states of cubic symmetry break up when this perturbation is applied.

\bigskip
\noindent\textbf{Solution.}

The cubic symmetry group is $\mathscr{O}$ (order 24) with 5 irreducible representations of dimensions 1, 1, 2, 3, 3 (denoted $\Gamma_1, \Gamma_2, \Gamma_3, \Gamma_4, \Gamma_5$ following the text's convention from the character table (10.91)).

The symmetry of a rectangular parallelepiped (with square top/bottom) is $D_4$ (order 8), which is a subgroup of $\mathscr{O}$. From Problem 16, $D_4$ has 5 irreducible representations of dimensions 1, 1, 1, 1, 2 (call them $\Gamma_1', \Gamma_2', \Gamma_3', \Gamma_4', \Gamma_5'$).

To find how levels split, we must determine how each irreducible representation of $\mathscr{O}$ decomposes when restricted to the subgroup $D_4$. This is done by comparing characters.

\medskip
\noindent\textbf{Character table of $\mathscr{O}$ (from Eq.~(10.91)):}

\[
\begin{array}{c|ccccc}
\mathscr{O} & \mathcal{C}_1 & \mathcal{C}_2 & \mathcal{C}_3 & \mathcal{C}_4 & \mathcal{C}_5 \\
& 1 & 6 & 3 & 8 & 6 \\
\hline
\Gamma_1 & 1 & 1 & 1 & 1 & 1 \\
\Gamma_2 & 1 & -1 & 1 & 1 & -1 \\
\Gamma_3 & 2 & 0 & 2 & -1 & 0 \\
\Gamma_4 & 3 & 1 & -1 & 0 & -1 \\
\Gamma_5 & 3 & -1 & -1 & 0 & 1
\end{array}
\]

We identify which elements of $\mathscr{O}$ belong to $D_4$. Taking $D_4$ as the subgroup preserving the $z$-axis (square cross-section in $xy$-plane), the 8 elements of $D_4$ within $\mathscr{O}$ are:
\begin{itemize}
\item From $\mathcal{C}_1$: $e$ (1 element)
\item From $\mathcal{C}_2$: $R_z(\pm 90^\circ)$ (2 elements)
\item From $\mathcal{C}_3$: $R_z(180^\circ)$ (1 element)
\item From $\mathcal{C}_5$: $R_x(180^\circ)$, $R_y(180^\circ)$ (2 elements from the face-diagonal class --- actually these are 180$^\circ$ about axes in the $xy$-plane)
\item From $\mathcal{C}_3$: $R_x(180^\circ), R_y(180^\circ)$ (2 more from edge-midpoint axes)
\end{itemize}

More carefully, the 8 elements of $D_4 \subset \mathscr{O}$ consist of:
\begin{itemize}
\item $e$
\item $C_4, C_4^3$ (90$^\circ$ and 270$^\circ$ about $z$) --- from $\mathcal{C}_2$ of $\mathscr{O}$
\item $C_4^2$ (180$^\circ$ about $z$) --- from $\mathcal{C}_3$
\item $C_2^x, C_2^y$ (180$^\circ$ about $x$, $y$ axes) --- from $\mathcal{C}_3$ of $\mathscr{O}$
\item $C_2^{d_1}, C_2^{d_2}$ (180$^\circ$ about face diagonals in $xy$-plane) --- from $\mathcal{C}_5$ of $\mathscr{O}$
\end{itemize}

Now we restrict each $\mathscr{O}$-representation to these 8 elements and decompose using the $D_4$ character table from Problem 16. The restricted characters on the 5 classes of $D_4$ are:

For $\Gamma_1$ of $\mathscr{O}$: $(1, 1, 1, 1, 1) \to \Gamma_1'$ of $D_4$.

For $\Gamma_2$: characters at $D_4$ elements: $e \to 1$, $C_4 \to -1$, $C_4^2 \to 1$, $C_2^{x,y} \to 1$, face diag $\to -1$. This is $\Gamma_3'$.

For $\Gamma_3$ (2D): $e \to 2$, $C_4 \to 0$, $C_4^2 \to 2$, $C_2^{x,y} \to 2$, face diag $\to 0$. Decomposing: $a_i = \frac{1}{8}\sum g_j \chi_j'^* \chi_j$.

Using $D_4$ character table: this gives $\Gamma_1' + \Gamma_3'$ (no splitting of 2D level --- it splits into $1+1$).

For $\Gamma_4$ (3D): $e \to 3$, $C_4 \to 1$, $C_4^2 \to -1$, $C_2^{x,y} \to -1$, face diag $\to -1$. Decomposing: $\Gamma_4' + \Gamma_5'$ (splits into $1 + 2$).

For $\Gamma_5$ (3D): $e \to 3$, $C_4 \to -1$, $C_4^2 \to -1$, $C_2^{x,y} \to -1$, face diag $\to 1$. Decomposing: $\Gamma_2' + \Gamma_5'$ (splits into $1 + 2$).

\medskip
\noindent\textbf{Summary of splittings:}

\[
\begin{array}{lcl}
\Gamma_1 \text{ of } \mathscr{O} \;(\text{1-fold}) &\to& \Gamma_1' \;(\text{no splitting}) \\
\Gamma_2 \text{ of } \mathscr{O} \;(\text{1-fold}) &\to& \Gamma_3' \;(\text{no splitting}) \\
\Gamma_3 \text{ of } \mathscr{O} \;(\text{2-fold}) &\to& \Gamma_1' + \Gamma_3' \;(\text{splits into } 1 + 1) \\
\Gamma_4 \text{ of } \mathscr{O} \;(\text{3-fold}) &\to& \Gamma_4' + \Gamma_5' \;(\text{splits into } 1 + 2) \\
\Gamma_5 \text{ of } \mathscr{O} \;(\text{3-fold}) &\to& \Gamma_2' + \Gamma_5' \;(\text{splits into } 1 + 2)
\end{array}
\]

The twofold degenerate level splits into two nondegenerate levels. Each threefold degenerate level splits into a nondegenerate level and a twofold degenerate level.

\end{document}
